\subsection{Definition of rational fractions}

DEFINITION 1. --- Let K be a commutative field and I a set. The field of fractions (I, p.~116) of the integral domain \(\mathrm{K}[(\mathrm{X}_{i})_{i\in I}]\) is denoted by \(\mathrm{K}((\mathrm{X}_{i})_{i\in I})\) or \(\mathrm{K}(\mathrm{X}_{i})_{i\in I}\) . Its elements are called rational fractions in the indeterminates \(X_{i}\) with coefficients in K.

For \(I = \{1, 2, \dots, n\}\) we write \(K(X_1, X_2, \dots, X_n)\) in place of \(K((X_i)_{i \in I})\) .

Let A be an integral domain and K its field of fractions. The ring \(\mathbf{A}[(\mathbf{X}_{i})_{i\in I}]\) may be identified with a subring of \(\mathbf{K}[(\mathbf{X}_{i})_{i\in I}]\) , hence also of \(\mathbf{K}((\mathbf{X}_{i})_{i\in I})\) . For each \(f\in\mathbf{K}[(\mathbf{X}_{i})_{i\in I}]\) there exists a non-zero element \(\alpha\) of A such that \(\alpha f\in\mathbf{A}[(\mathbf{X}_{i})_{i\in I}]\) . Hence every element of \(\mathbf{K}((\mathbf{X}_{i})_{i\in I})\) can be written as u/v where \(u,v\in\mathbf{A}[(\mathbf{X}_{i})_{i\in I}],v\neq0\) . Thus \(\mathbf{K}((\mathbf{X}_{i})_{i\in I})\) may be identified with the field of fractions of \(\mathbf{A}[(\mathbf{X}_{i})_{i\in I}]\) .

Now let K be a commutative field, I a set and \(J \subset I\) . Put \(\mathbf{B} = \mathbf{K}[(\mathbf{X}_i)_{i \in J}]\) , then \(\mathbf{K}[(\mathbf{X}_i)_{i \in I}] = \mathbf{B}[(\mathbf{X}_i)_{i \in I-J}]\) , and by what has been said, \(\mathbf{K}((\mathbf{X}_i)_{i \in I})\) may be identified with \(\mathbf{K}'((\mathbf{X}_i)_{i \in I-J})\) , where \(\mathbf{K}' = \mathbf{K}((\mathbf{X}_i)_{i \in J})\) .

